\documentclass[11pt,a4paper]{article}
\usepackage[T1]{fontenc}
\usepackage[utf8]{inputenc}
\usepackage{lmodern}
\usepackage[indonesian]{babel}
\usepackage[a4paper,margin=2.2cm]{geometry}
\usepackage{microtype}
\usepackage{amsmath,amssymb}
\usepackage{booktabs,tabularx,array}
\usepackage{enumitem}
\usepackage[most]{tcolorbox}
\usepackage{xcolor}
\usepackage{fancyhdr}
\usepackage{needspace}
\usepackage[hidelinks]{hyperref}
\setlength{\parindent}{0pt}
\setlength{\parskip}{5pt}
\setlist{itemsep=2pt,topsep=3pt}
\renewcommand{\arraystretch}{1.25}
\definecolor{biru}{HTML}{1F4E79}
\definecolor{hijau}{HTML}{2E7D32}
\definecolor{merah}{HTML}{B71C1C}
\definecolor{kuning}{HTML}{8D6E00}
\newtcolorbox{intinya}{colback=biru!6,colframe=biru,title=Intinya,fonttitle=\bfseries,boxrule=0.6pt,arc=2pt}
\newtcolorbox{hafal}{colback=hijau!6,colframe=hijau,title=Wajib hafal,fonttitle=\bfseries,boxrule=0.6pt,arc=2pt}
\newtcolorbox{awas}{colback=merah!5,colframe=merah,title=Jangan salah ucap,fonttitle=\bfseries,boxrule=0.6pt,arc=2pt}
\newtcolorbox{contoh}[1][]{colback=kuning!5,colframe=kuning,title={Contoh hitung #1},fonttitle=\bfseries,boxrule=0.6pt,arc=2pt}
\newcounter{tanya}
\newenvironment{tanya}[1]{\refstepcounter{tanya}\par\medskip\needspace{4\baselineskip}\noindent\textbf{\color{biru}T\thetanya.\ #1}\par\noindent\textbf{Jawab:}\ }{\par}
\pagestyle{fancy}\fancyhf{}
\fancyhead[L]{\small Bahan Belajar Proposal --- Abd. Rouf (223140118)}
\fancyhead[R]{\small\thepage}
\renewcommand{\headrulewidth}{0.4pt}
\newcommand{\judul}[2]{\begin{center}{\LARGE\bfseries\color{biru}#1}\par\vspace{4pt}{\large #2}\par\vspace{2pt}{\small Pegangan pribadi untuk bimbingan dan ujian proposal. Isi mengikuti \textit{Proposal\_Skripsi\_ABD\_ROUF.docx}.}\end{center}\vspace{-4pt}\hrule\vspace{6pt}}

\begin{document}
\judul{Bab I --- Pendahuluan}{Analisis QoS Jaringan Berbasis Crowdsourcing pada Aplikasi Presensi Menggunakan K-Means}

\begin{intinya}
Karyawan PT Semeru Agung Mandiri sering mengeluh presensi lambat atau gagal terkirim. Aplikasi sudah \textit{offline-first}, jadi presensi tidak hilang. Tetapi itu hanya obat gejala. \textbf{Penyebabnya (kualitas jaringan di lokasi kerja) belum pernah diukur.}
Penelitian ini memakai aplikasi presensi itu sendiri sebagai \textbf{alat ukur jaringan}: setiap kali karyawan absen, HP-nya mengukur kualitas jaringan. Hasilnya dinilai dengan standar \textbf{TIPHON}, perbedaannya diuji dengan \textbf{Kruskal-Wallis}, dan polanya dikelompokkan dengan \textbf{K-Means}.
\end{intinya}

\section*{1. Alur logika Bab I (hafalkan urutannya)}
Bab I adalah ``cerita'' yang membawa penguji dari masalah ke solusi. Urutannya:
\begin{enumerate}[label=\textbf{\arabic*.}]
  \item \textbf{Aplikasi presensi bergantung pada jaringan.} Data, foto, dan lokasi dikirim lewat Wi-Fi/seluler. Penelitian lain menambah mode offline karena koneksi lapangan tidak stabil.
  \item \textbf{Di perusahaan ini terjadi hal yang sama.} Ada keluhan; mode offline hanya mengatasi gejala, penyebabnya tidak diukur.
  \item \textbf{Bukti awal (studi pendahuluan)} dari data presensi 20 April--27 Mei 2026 (Tabel 1.1).
  \item \textbf{Inti masalah:} tidak ada \textit{data objektif} tentang kualitas jaringan saat presensi --- di lokasi mana, operator apa, seberapa buruk.
  \item \textbf{Cara mengukur yang objektif} = parameter QoS (delay, jitter, packet loss, throughput) + kategori TIPHON.
  \item \textbf{Kelemahan cara lama:} penelitian QoS biasanya manual di satu titik dengan Wireshark atau nPerf, padahal kualitas berubah menurut waktu (contoh jam sibuk: packet loss 26,5\%, delay 485 ms). Aplikasi speedtest juga mengukur ke server mereka sendiri, bukan server perusahaan.
  \item \textbf{Membandingkan rata-rata saja tidak cukup:} data jaringan tidak normal, jadi perlu uji statistik (Kruskal-Wallis, contoh Makarim 2026).
  \item \textbf{Solusi pengukuran:} \textit{crowdsourcing} --- diukur otomatis di HP banyak pengguna (Wamser 2021; Wu 2017, 5 juta pengukuran dari 2.351 HP).
  \item \textbf{Solusi pengelompokan:} K-Means, karena indeks TIPHON adalah rata-rata sehingga ``lambat tapi stabil'' bisa sama nilainya dengan ``cepat tapi sering gagal''.
  \item \textbf{Tiga celah penelitian} (lihat bagian 3).
  \item \textbf{Mengapa perlu:} presensi = dasar gaji; perbaikan jaringan butuh biaya sehingga harus tepat sasaran; aplikasi sudah ada di semua lokasi sehingga murah.
  \item \textbf{Kesimpulan:} tambahkan modul telemetri $\rightarrow$ data 4 minggu $\rightarrow$ TIPHON $\rightarrow$ Kruskal-Wallis $\rightarrow$ K-Means $\rightarrow$ profil lokasi.
\end{enumerate}

\section*{2. Angka studi pendahuluan (Tabel 1.1)}
\begin{hafal}
\begin{tabularx}{\linewidth}{@{}lX@{}}
Periode & 20 April -- 27 Mei 2026 (38 hari) \\
Jumlah data & 1.911 catatan presensi, 88 pengguna \\
Lokasi & 14 dari 16 lokasi terdaftar dipakai \\
Jam padat & 91,2\% presensi masuk pukul 07.00--08.59 \\
Wi-Fi kantor & hanya 4,3\% $\Rightarrow$ kualitas ditentukan operator seluler \\
Mode GAMAS & 282 catatan (14,8\%) \\
Data tidak seimbang & 6 dari 14 lokasi $<$ 30 catatan; 4 lokasi terbesar = 68,9\% \\
Data jaringan & tidak ada; penanda offline bernilai 0 di semua catatan \\
\end{tabularx}
\end{hafal}
Arti setiap temuan (kolom ``Dampak''):
\begin{itemize}
  \item \textbf{Keluhan lambat/gagal} $\rightarrow$ rekap gaji terlambat, HRD cek manual.
  \item \textbf{Wi-Fi cuma 4,3\%} $\rightarrow$ yang menentukan kualitas adalah \textit{BTS operator}. Ini yang membuat topik ini ``jaringan/infrastruktur fisik''.
  \item \textbf{91,2\% pagi} $\rightarrow$ gangguan jam itu kena banyak orang sekaligus.
  \item \textbf{Tidak ada data jaringan} $\rightarrow$ keluhan tidak bisa dibuktikan. \textit{Inilah inti masalah.}
  \item \textbf{Data tidak seimbang} $\rightarrow$ lokasi kecil perlu perlakuan khusus (aturan minimal 30 sesi).
  \item \textbf{Ukur manual mahal} $\rightarrow$ ukur sesaat di satu titik tidak mewakili 16 lokasi.
\end{itemize}

\section*{3. Tiga celah penelitian (\textit{research gap})}
\begin{enumerate}
  \item Penelitian QoS mengukur \textbf{manual, di satu titik, pada waktu tertentu} (Rosid 2023; Fiqria \& Purnama 2024).
  \item Penelitian K-Means kualitas jaringan memakai \textbf{data pihak ketiga} (Ookla) atau \textbf{alat pemantau} (Cacti), bukan data dari aplikasi yang dipakai sehari-hari (Irwansyah 2025; Mubarak 2023).
  \item Penelitian aplikasi presensi hanya menambah \textbf{mode offline tanpa mengukur jaringannya} (Indah \& Wati 2025; Rokhman 2026).
\end{enumerate}
\textbf{Kebaruan} = gabungan ketiganya: aplikasi presensi sebagai instrumen QoS \textit{crowdsourcing}, diuji beda antarlokasi \& antaroperator, lalu dikelompokkan K-Means.

\section*{4. Batasan, rumusan, tujuan, manfaat}
\textbf{Batasan (6 butir)} --- cara mudah mengingat: \textbf{A-U-J-P-D-A}
\begin{enumerate}[label=\textbf{\Alph*}]
  \item[\textbf{A}] \textbf{Aplikasi}: Android versi 2.0.1; iOS tidak.
  \item[\textbf{U}] \textbf{Unit}: sesi presensi yang lolos \textit{geofence}; sinkronisasi latar (WorkManager) tidak dicatat.
  \item[\textbf{J}] \textbf{Jalur}: HTTPS HP $\rightarrow$ server perusahaan; \texttt{generate\_204} Google sebagai pembanding.
  \item[\textbf{P}] \textbf{Parameter}: delay (RTT tingkat aplikasi), jitter, tingkat kegagalan (pengganti packet loss), throughput unggah, lama putus; plus jenis jaringan \& operator. \textit{Sinyal dan lokasi BTS tidak diukur.}
  \item[\textbf{D}] \textbf{Data}: 4 minggu, pseudonim, tanpa nama/email/GPS/foto; uji beda hanya lokasi $\geq$ 30 sesi.
  \item[\textbf{A}] \textbf{Analisis}: TIPHON, Kruskal-Wallis, K-Means; validasi di meja (PCAPdroid, Wireshark). Tidak memperbaiki jaringan, hanya rekomendasi.
\end{enumerate}

\textbf{Rumusan masalah} dan \textbf{tujuan} saling berpasangan:
\begin{center}
\begin{tabularx}{\linewidth}{@{}p{0.06\linewidth}XX@{}}
\toprule
No. & Rumusan (bertanya) & Tujuan (menjawab) \\
\midrule
1 & Bagaimana kualitas jaringan per lokasi \& operator menurut TIPHON, dan apakah berbeda signifikan? & Mendeskripsikan dengan TIPHON + menguji beda dengan Kruskal-Wallis \\
2 & Pola apa yang terbentuk dengan K-Means, dan lokasi mana didominasi pola buruk? & Mengelompokkan sesi (delay, jitter, gagal) + menyusun profil lokasi \\
\bottomrule
\end{tabularx}
\end{center}
\textbf{Manfaat:} teoretis (cara ukur QoS \textit{crowdsourcing} dengan aplikasi operasional + penerapan Kruskal-Wallis \& K-Means); praktis untuk \textbf{perusahaan} (penempatan AP/pilih operator), \textbf{HRD} (telusuri keluhan), \textbf{pengembang} (dua alamat membedakan masalah server vs jaringan).

\section*{5. Kemungkinan pertanyaan penguji}
\setcounter{tanya}{0}
\begin{tanya}{Ini skripsi jaringan atau skripsi aplikasi?}
Jaringan. Yang dianalisis adalah \textbf{kinerja jaringan} (QoS), bukan fitur aplikasi. Aplikasi hanya \textbf{alat ukur}, sama seperti Wireshark di penelitian QoS lain. Bidang minat: AI \& Networks, topik \textit{Network Management}.
\end{tanya}
\begin{tanya}{Di mana unsur perangkat kerasnya?}
Yang diukur adalah infrastruktur fisik: \textbf{BTS operator} dan \textbf{access point Wi-Fi kantor}. HP karyawan berperan sebagai sensor. Karena hanya 4,3\% presensi lewat Wi-Fi kantor dan sebagian besar sisanya lewat data seluler, kualitasnya banyak ditentukan BTS di sekitar lokasi.
\end{tanya}
\begin{tanya}{Kenapa tidak datang ke tiap lokasi dan pakai Wireshark saja?}
Ukur manual hanya memotret satu titik pada satu waktu, mahal untuk 16 lokasi, dan tidak menangkap jam padat pagi. \textit{Crowdsourcing} mengukur tiap hari, di semua lokasi, pada saat karyawan benar-benar absen. Wireshark tetap dipakai --- untuk \textbf{memvalidasi} alat ukur.
\end{tanya}
\begin{tanya}{Apa bedanya dengan Speedtest/nPerf?}
Speedtest mengukur ke server miliknya sendiri dan harus dijalankan manual. Penelitian ini mengukur ke \textbf{server perusahaan yang dipakai presensi}, otomatis, dan ada alamat pembanding untuk memisahkan masalah server vs jaringan.
\end{tanya}
\begin{tanya}{Apa kebaruannya?}
Belum ada penelitian yang memakai aplikasi presensi sebagai instrumen QoS \textit{crowdsourcing}, menguji beda antarlokasi \& antaroperator, lalu mengelompokkan dengan K-Means. Sebutkan tiga celah di bagian 3.
\end{tanya}
\begin{tanya}{Kenapa tidak cukup TIPHON saja? Kenapa perlu K-Means?}
Indeks TIPHON adalah rata-rata empat parameter. Koneksi lambat tapi stabil bisa mendapat indeks sama dengan koneksi cepat tapi sering gagal. K-Means memisahkan pola berdasarkan semua parameter sekaligus.
\end{tanya}
\begin{tanya}{Kenapa perlu uji Kruskal-Wallis? Bukankah rata-rata sudah kelihatan beda?}
Selisih angka belum tentu perbedaan nyata. Makarim (2026): ping beda signifikan, tapi unggah tidak (p = 0,975) walaupun angkanya beda. Data jaringan tidak normal, jadi dipakai uji nonparametrik.
\end{tanya}
\begin{tanya}{Kenapa kekuatan sinyal dan lokasi BTS tidak diukur?}
Butuh izin lokasi presisi/latar dan data BTS dari operator; menambah risiko privasi dan baterai. Ini ditulis sebagai keterbatasan. Operator dan jenis jaringan tetap dicatat sebagai konteks.
\end{tanya}
\begin{tanya}{Apa tidak melanggar privasi karyawan?}
Ada izin tertulis perusahaan, pemberitahuan di aplikasi, identitas diganti pseudonim, tanpa nama/email/GPS/foto, hanya admin yang bisa melihat, dan hasil disajikan agregat (UU PDP 27/2022).
\end{tanya}
\begin{tanya}{Kenapa hanya 4 minggu?}
Cukup untuk menangkap pola hari kerja berulang (masuk \& pulang, hari berbeda) dengan perkiraan $\pm$1.400 catatan / maks. $\pm$2.800 sesi. Lokasi kecil dilaporkan deskriptif.
\end{tanya}

\begin{awas}
\begin{itemize}
  \item Jangan bilang ``data \textbf{anonim}'' --- yang benar \textbf{pseudonim}.
  \item Jangan bilang ``mengukur \textbf{packet loss}'' --- yang benar ``tingkat kegagalan sebagai \textbf{pengganti} packet loss''.
  \item Jangan bilang ``penelitian R\&D / membuat aplikasi'' --- yang benar \textbf{kuantitatif deskriptif-komparatif}; modul = instrumen.
  \item Jangan bilang ``ping'' --- yang diukur \textbf{RTT tingkat aplikasi} (HTTP), bukan ICMP.
  \item Penelitian \textbf{tidak memperbaiki} jaringan, hanya memberi rekomendasi.
\end{itemize}
\end{awas}

\section*{6. Latihan mandiri (jawab tanpa melihat)}
\begin{enumerate}
  \item Sebutkan inti masalah dalam satu kalimat.
  \item Sebutkan 3 angka studi pendahuluan yang paling kuat dan artinya.
  \item Sebutkan 3 celah penelitian beserta 1 sumber masing-masing.
  \item Jelaskan dengan bahasa sendiri kenapa TIPHON saja tidak cukup.
  \item Sebutkan 6 batasan dengan singkatan A-U-J-P-D-A.
\end{enumerate}
\end{document}
